\begin{abstract}
Tool-using language-model agents turn untrusted text into structured calls that
commit external state.  A single call can create an object, disclose data to
several principals, and change visibility.  A monitor that sees only the tool
name, call, or raw arguments may therefore merge state changes that require
different authorization decisions.  We study \emph{tool-effect binding}: the
representation needed to connect a proposed call to the security-relevant
effects it will commit.  We show that when a representation merges effects
that an admissible policy distinguishes, every monitor over that representation
must either allow an unauthorized effect or withhold authorized work.

We introduce counterfactual registration, which treats a tool descriptor as a
hypothesis, executes controlled call and state variants, and refines the
descriptor when the observed effects expose a missing distinction.  Source
execution determines what changed; a separate policy oracle determines whether
the change must be authorized independently.  A common-field representation
leaves 118 authorization-separating pairs across 56 source-executed calls,
whereas typed effect atoms remove every observed collision.  The same
representation removes all observed collisions in 32 held-out ToolSandbox
contexts and exactly realizes their frozen 232-query policy relation.

Finally, we use a conservative registered descriptor in a deterministic
provenance-origin monitor.  On 608 scope-aligned AgentDojo attacks under
DeepSeek, attack success falls from 1.0\% to 0.0\%, and every executed call has
a matching pre-commit decision.  A same-model prompting defense reaches the
same attack rate, while benign non-inferiority is not established.  We
therefore treat this runtime experiment as a bounded mediation case study; the
main representation claim rests on source-executed collisions and the frozen
finite-policy test.
\end{abstract}
